\section{Preliminaries}
\label{sec:preliminaries}

\subsection{Compact cores and auxiliary caps}

\begin{lemma}\label{lem:compact-core}
The sufficiently large coordinate spheres bound connected compact
truncations of $\Omega$. Each component of $S$ can be capped on its
other side to produce a smooth complete connected one-ended manifold
$X$ without boundary, with the orientation extending that of $\Omega$.
Let $O$ denote the compact cap set, including its
boundary. For competitors containing $O$, perimeter depends only on
the metric on $\Omega$ and its boundary trace.
\end{lemma}
\begin{proof}
The end metric is uniformly comparable to $b$ sufficiently far out.
Radial infinity escapes compact subsets: otherwise a sequence tending
to radial infinity would accumulate in a compact coordinate
neighborhood, while pairwise disjoint background balls of a fixed small
radius about a subsequence would have uniformly positive metric
volumes and lie in a fixed compact neighborhood. This is impossible.
A large coordinate sphere separates the product end from an inner
side. If the latter were noncompact, a smooth compact exhaustion would
yield an end separated from the given product end. One may choose this
exhaustion with finitely many complement components at each stage,
since its smooth compact frontier has finitely many components.
Nested extraction then contradicts the one-end assumption. The inner
side is therefore compact. Paths between its points can be pushed
back across the product end, so the truncations can be taken connected.

Every component of $S$ is a closed orientable surface and bounds a
handlebody. Orient each handlebody and glue by a map reversing the
induced boundary orientations.  The collars then extend the orientation
of $\Omega$ across each seam. Extend the metric smoothly across $S$.
The added region is compact; completeness
at infinity is unchanged. For a set containing $O$ up to a null set,
its characteristic function is constant almost everywhere in the
cap interiors. Its perimeter therefore uses no metric values there.
\end{proof}

We use the ordinary perimeter of a finite-perimeter set in the capped
manifold. In particular it is not a relative perimeter that discards
the original boundary.

\begin{lemma}\label{lem:cut-enclosure}
Every admissible cut from Definition~\ref{def:minimum-area} determines
a precompact finite-perimeter enclosure of $O$ whose perimeter is at
most the area of that cut. Conversely, a precompact enclosure of $O$
with smooth free boundary in $\operatorname{int}\Omega$ and connected
exterior gives an admissible cut with the same area. If this enclosure
is outward minimizing, its boundary realizes the minimum enclosing
area of its own closed exterior.
\end{lemma}
\begin{proof}
For an admissible $D$, complement its manifold interior in $X$.
The result contains $O$ and is precompact because $D$ contains the
whole sufficiently distant end. Local boundary coordinates show that
its perimeter is bounded by the entire intrinsic area of $\partial D$,
including any portions on $S$. This also follows by smoothing corners
and passing to the perimeter lower limit. For the converse, take the
closed exterior of the smooth enclosure. It is connected, has the
required end, and has precisely the stated intrinsic boundary.
Any admissible cut of this new exterior, complemented in $X$, encloses
the original enclosure; outward minimization gives the final claim.
\end{proof}

The compactness, lower semicontinuity, coarea, slicing, reduced-boundary
structure, and union/intersection inequalities for finite-perimeter
sets will be used below. These are local facts and hold in smooth
Riemannian coordinate charts; see \citet[Chapters 12--18]{Maggi2012}.
They also give compactness and lower semicontinuity for a uniformly
convergent family of uniformly positive continuous metrics on a fixed
compact set, by comparison with smooth metrics. A free locally
perimeter-minimizing boundary is smooth in ambient dimension three;
one may use the smooth-ambient regularity theorem in
\citet[Chapter 7, Section 5, Theorem 5.8]{Simon2014}, followed by
elliptic bootstrapping. We invoke this regularity only after excluding
contact with every obstacle and outer truncation.

\subsection{Mass covariance and small metric changes}

\begin{proposition}\label{prop:mass-covariance}
Suppose $g-b\in C^{2,\alpha}_\tau$, $\tau>3/2$, and
$\int V_0|R_g+6|\dd V_g<\infty$ on the end. The fluxes
\eqref{eq:ah-mass} transform as a Lorentz covector under hyperbolic
isometries of the asymptotic chart. Such chart changes preserve the
weighted decay and integrability hypotheses used here. Consequently
a future-timelike original covector admits a chart with
$p_0=m_{\AH}$ and $p_1=p_2=p_3=0$.
\end{proposition}
\begin{proof}
The potentials in \eqref{eq:background} are the ambient coordinate
functions of the future unit hyperboloid. Hyperbolic isometries act on
their span by Lorentz transformations. The flux integrand is tensorial
in $b$, linear in $V$, and natural under these isometries.
It remains to justify replacing the transformed integration surfaces
by the original large coordinate spheres. The static identity
$\nabla_b^2V=Vb$ gives
\[
 \divg_b\mathbb U(V,e)=V\,D R_b(e).
\]
The difference between $D R_b(e)$ and $R_g+6$ is bounded by a constant
times
\[
 |e|_b^2+|\nabla_be|_b^2+|e|_b|\nabla_b^2e|_b.
\]
Its $V_0$-weighted integral over $r>D$ is
$O(D^{3-2\tau})$, since $\dd V_b$ is comparable to
$r\dd r\dd A_\sigma$. Also $|V_i|\le V_0$.
Thus the divergence is absolutely integrable on the end for each
potential. The divergence theorem on regions between the two
families of spheres makes their flux difference tend to zero.
This is the covariance mechanism of
\citet[Proposition 2.2, Theorem 2.3, and Section 3]{CH2003} in the
normalization \eqref{eq:ah-mass}.

A fixed hyperbolic isometry changes distance from the chosen origin
by a bounded amount. The exponential weights and $V_0$ are therefore
comparable before and after the change, and uniform background balls
are carried to uniform background balls. The rest-chart assertion is
the elementary Lorentz normal form of a future-timelike covector.
\end{proof}

The boundary sign and the enclosing-area definition are intrinsic.
We will apply Proposition~\ref{prop:mass-covariance} to the original
data before constructing any comparison metric. No assertion that
intermediate covectors remain timelike is required.

If $(1-\eps)g\le\widetilde g\le(1+\eps)g$ as quadratic forms,
then for every two-surface its area changes by factors between
$1-\eps$ and $1+\eps$. Taking infima over the same admissible domains
gives the same comparison for $A_{\min}$. This observation will be
used when passing through the preliminary conformal approximation
and the later corner smoothing.

\subsection{The AF comparison theorem}

For an AF metric $h$ in a Euclidean end chart $y$, our normalization is
\begin{equation}\label{eq:adm-mass}
 m_{\AF}(h)=\frac1{16\pi}\lim_{R\to\infty}
 \int_{|y|=R}(\partial_jh_{ij}-\partial_ih_{jj})
                  n^i\dd A_{\mathrm{Eucl}}.
\end{equation}

\begin{theorem}[AF Riemannian Penrose inequality]\label{thm:af-penrose}
Let $(Y,h)$ be a complete connected one-ended three-dimensional AF
manifold with smooth compact outer-minimizing minimal boundary.
Assume $R_h\ge0$, $h-\delta=O_2(|y|^{-\beta})$ for some
$\beta>1/2$, and $R_h=O(|y|^{-q})$ for some $q>3$.
If the total boundary area is $A$, then
\begin{equation}\label{eq:af-penrose}
 m_{\AF}(h)\ge\sqrt{A/(16\pi)}.
\end{equation}
Disconnected boundary is allowed.
\end{theorem}

This is the dimension-three case of the published
\citet[Theorem 1.4]{BrayLee2009}; its asymptotic definition is on
page 82 and its boundary convention on page 83.
Only the inequality is used. The smoothness, scalar sign, and
outer-minimization required here will be established for the AF
comparison exterior, rather than assumed for either original AH
metric.

\subsection{Local scalar elliptic estimates}

We specify the standard local analytic inputs, separating them from
the coupled transmission estimates proved in
Section~\ref{sec:transmission-analysis}.

\begin{proposition}[Scalar local estimates]\label{prop:elliptic-inputs}
On fixed coordinate balls in dimension three the following estimates
hold, with constants depending on the indicated ellipticity,
coefficient, and geometry bounds.
\begin{enumerate}[label=\textup{(\roman*)}]
\item Bounded measurable uniformly elliptic divergence matrices admit
local subsolution bounds, weak Harnack estimates for nonnegative
supersolutions, and inhomogeneous Harnack and H\"older estimates for
nonnegative solutions. Bounded divergence drifts multiplying the
unknown, vector forcing in $L^q$, $q>3$, and scalar forcing in $L^p$,
$p>3/2$, are permitted in the corresponding estimates.
\item For Laplace equations with uniformly positive Lipschitz metric
coefficients in interior or flattened smooth boundary charts,
divergence forcing in $L^q$ gives local $W^{1,q}$ estimates.
Scalar forcing in $L^2$ can be included for $3<q<6$.
The same statements hold with local zero Dirichlet data or natural
Neumann data, with the divergence forcing included in the conormal
current. A boundary datum in $L^{p_1}$, $p_1>2$, permits a choice
$3<q<3p_1/2$ sufficiently close to $3$.
\item Smooth-coefficient scalar Dirichlet and conormal problems have
the usual Schauder estimates, including divergence-form
$C^{1,\gamma}$ estimates when the uniformly elliptic coefficients
and vector forcing are $C^\gamma$, with boundary data of the
corresponding regularity. These estimates are used only after their
coefficient hypotheses have been obtained.
\item For a bounded measurable symmetric uniformly elliptic matrix
$a$ on a ball, its positive Dirichlet Green kernel is bounded above
by $C|z-z'|^{-1}$, and below by $c|z-z'|^{-1}$ when the pole and
the nearby evaluation point lie in a fixed smaller ball.
\end{enumerate}
\end{proposition}

The local boundedness, Harnack, and H\"older estimates are the scalar
divergence theory of \citet{Serrin1964} and
\citet[Chapter 8]{GT2001}. The Laplace divergence estimates are the
local Calder\'on--Zygmund estimates, with smooth-boundary
localization; see \citet{Stein1970,GT2001}.
For natural Neumann data, an explicit estimate is given by
\citet[Appendix, proof of Lemma 6.1, (6.2), p.\ 19
of arXiv v3]{ChoiKim2013} on bounded $C^1$ localizations.
Our Lipschitz metric coefficients are VMO. A cutoff retains the
full conormal current and the larger-patch $W^{1,2}$ norm.
In dimension three the scalar and boundary forcing exponents are
$3q/(q+3)$ and $2q/3$, respectively, so approximation permits
$3<q<\min\{6,3p_1/2\}$ for the data in part~(ii).
For the boundary exponent in part~(ii), the trace embedding of
$W^{1,q'}$ into $L^{2q/(2q-3)}$ on a two-dimensional face allows
pairing with $L^{p_1}$ whenever $q<3p_1/2$.
Scalar $L^2$ forcing belongs locally to $W^{-1,q}$ when $q<6$.
Part~(iv) is the local Green comparison of
\citet[Theorem 7.1 and the discussion on pp.\ 66--67]{LSW1963}.

Natural homogeneous flux permits even reflection in a flattened
boundary chart, with the full current reflected. Smooth prescribed
flux errors can be placed in a bounded smooth vector field.
Fermi-coordinate metric matrices reflected at a smooth face are
Lipschitz. Constant Dirichlet data for a scalar height permit odd
reflection after subtraction of that constant. We spell out each
reflection where it is needed; none of these scalar statements is a
general nonlinear transmission theorem.

For the electrostatic equation we also use an exact degenerate
gradient regularity statement. Suppose
\[
 A(z,0)=0,\qquad
 \lambda|\xi|I\le A_\xi(z,\xi)\le C|\xi|I,
 \qquad
 |A(z,\xi)-A(z',\xi)|\le C|z-z'|\,|\xi|^2,
\]
with symmetric $A_\xi$, continuous dependence, and smooth dependence
on $\xi\ne0$. Then a bounded weak $W^{1,3}$ solution of
$\partial_iA^i(z,\dd u)=0$ is locally $C^{1,\gamma}$ for some
$\gamma>0$, with uniform estimates from the structural constants
and a local energy bound. Indeed $A_\xi(z,0)=0$ extends continuously,
and integration along a radial segment gives $|A(z,\xi)|\le C|\xi|^2$.
The assumptions of \citet[Theorem 1.1, (1.5)--(1.8)]{AraujoZhang2020}
therefore hold with $n=p=3$, zero forcing in $L^\infty$, and spatial
H\"older exponent $\sigma=1/2$.  We use only an exponent
$0<\gamma<\min\{\alpha_m,1/4\}$, where $\alpha_m>0$ is a
constant-coefficient gradient exponent for the same structural constants.
No optimal exponent or endpoint estimate is required.
If these inequalities initially hold only on a known
bounded gradient range, the constitutive law must first be extended
with the same structure; that extension is constructed in
Section~\ref{sec:stress-extension}.

All global invertibility, weighted end estimates, seam regularity,
continuation, and uniform large-parameter conclusions used later are
proved below. The local estimates in this subsection supply only
their explicitly stated scalar ingredients.
